
\begin{exercise}[Find the mistake \Coffeecup]
    Let \((X_n)_{n\in \nat}\) and \(X\) be random variables in \(\real^\dims\).
    Which step in the following chain of implications is \textbf{wrong}?
    \[
        X_n \overset{d}\to X
        \overset{?}\implies X_n - X \overset{d}\to 0
        \overset{?}\implies X_n - X \overset{p}\to 0
        \overset{?}\implies X_n \overset{p}\to X.
    \]
    Justify all other steps and find a counterexample for the wrong step.
    \begin{solution}
        The first step is wrong. Consider \(X_n = -X\) with
        \(\prob{X=-1}=\prob{X=1}=\frac12\). Then
        as in  Exercise \ref{ex: convergence in distribution but not in
        probability} \(X_n \to X\), but \(X_n - X = -2X\) does not converge to
        zero in distribution.

        All the other steps are correct. The second step is the special case
        that convergence in distribution to a constant implies convergence in
        probability to that constant (Theorem \ref{thm: convergence
        implications}~\ref{it: d to p}). The third step follows from 
        \[
            \prob{\norm{X_n - X} > \epsilon}
            =\prob[\big]{\norm{(X_n - X) - 0} > \epsilon}
            \to 0.
            \qedhere
        \]
    \end{solution}
\end{exercise}
